% =====================================================================
%  Print design for the C64 Commander / C64U Remote manuals.
%
%  The reference is the 1982 Commodore 64 User's Guide: a friendly book
%  a reader settles into, not a datasheet. What that means in practice:
%  a warm oldstyle body face at a comfortable measure, a grotesque sans
%  for anything structural, one accent hue per chapter carried through
%  the opener, the running head and the index, and enough air that a
%  page never looks like work.
%
%  Body:     TeX Gyre Pagella (Palatino) - warm, generous, made for books
%  Headings: TeX Gyre Heros (Helvetica)  - the era's own structural sans
%  Code:     TeX Gyre Cursor (Courier)   - what a C64 listing looks like
% =====================================================================
\usepackage[T1]{fontenc}
\usepackage[utf8]{inputenc}
\usepackage{tgpagella}
\renewcommand{\sfdefault}{qhv}
\renewcommand{\ttdefault}{qcr}
\usepackage[activate={true,nocompatibility},final,tracking=true,factor=1100,stretch=12,shrink=12]{microtype}
\usepackage{xcolor}
\usepackage{graphicx}
\usepackage{booktabs}
\usepackage{longtable}
\usepackage{array}
\usepackage{ragged2e}
\usepackage{enumitem}
\usepackage{tcolorbox}
% `calc` also gives pandoc's longtable column specs their \real macro.
\usepackage{calc}
\usepackage{eso-pic}
\usepackage[hidelinks]{hyperref}
\tcbuselibrary{skins,breakable}

% ---------------------------------------------------------------- page
% 210 - 30 - 45 leaves a 135 mm measure: about 70 characters a line at
% 11 pt, which is the range long-form text is set in. The inner margin
% is the narrower of the two because the gutter of a bound book eats
% part of it.
\setstocksize{297mm}{210mm}
\settrimmedsize{297mm}{210mm}{*}
\settypeblocksize{225mm}{135mm}{*}
\setlrmargins{30mm}{*}{*}
\setulmargins{32mm}{*}{*}
\setheadfoot{14pt}{26pt}
\setheaderspaces{*}{16pt}{*}
\checkandfixthelayout
\setlength{\parindent}{0pt}
\setlength{\parskip}{0.62\baselineskip}
\linespread{1.06}
\raggedbottom
\sloppybottom
\midsloppy

% -------------------------------------------------------------- colour
% One hue per chapter, in the order the chapters run. Warm enough to sit
% on a cream page, distinct enough to tell chapters apart at a glance.
\definecolor{c64ink}{HTML}{1C1B19}
\definecolor{c64muted}{HTML}{5C5952}
\definecolor{c64rule}{HTML}{D8D2C6}
\definecolor{c64paper}{HTML}{FBF9F4}
\definecolor{chapA}{HTML}{2A6F97}
\definecolor{chapB}{HTML}{1F7A6D}
\definecolor{chapC}{HTML}{6D597A}
\definecolor{chapD}{HTML}{B5651D}
\definecolor{chapE}{HTML}{8A6A3B}
\definecolor{chapF}{HTML}{40655E}
\definecolor{chapG}{HTML}{9C4722}
\definecolor{chapH}{HTML}{A4243B}
\definecolor{chapI}{HTML}{34568B}
\colorlet{accent}{chapA}

% Rotates the accent as each chapter opens, so the opener, the headings,
% the figure labels and the running head of one chapter always agree.
\newcounter{accentindex}
\newcommand{\rotateaccent}{%
  \stepcounter{accentindex}%
  \ifcase\numexpr\value{accentindex}-1\relax
    \colorlet{accent}{chapA}\or\colorlet{accent}{chapB}\or\colorlet{accent}{chapC}%
    \or\colorlet{accent}{chapD}\or\colorlet{accent}{chapE}\or\colorlet{accent}{chapF}%
    \or\colorlet{accent}{chapG}\or\colorlet{accent}{chapH}\or\colorlet{accent}{chapI}%
  \else\colorlet{accent}{chapA}\fi}

\pagecolor{c64paper}
\color{c64ink}

% ------------------------------------------------------------ headings
% A chapter opens on a fresh recto with the number set large and pale in
% the chapter's own hue, the title beneath it, and a rule the width of
% the measure. The eye lands on the title, then the rule sends it into
% the text.
\makeatletter
\makechapterstyle{c64}{%
  \setlength{\beforechapskip}{0pt}
  \setlength{\midchapskip}{6pt}
  \setlength{\afterchapskip}{26pt}
  \renewcommand{\chapnamefont}{\normalfont}
  \renewcommand{\chapnumfont}{\sffamily\bfseries\fontsize{54}{54}\selectfont\color{accent}}
  \renewcommand{\chaptitlefont}{\sffamily\bfseries\fontsize{25}{29}\selectfont\color{c64ink}}
  \renewcommand{\printchaptername}{}
  \renewcommand{\chapternamenum}{}
  \renewcommand{\printchapternum}{\rotateaccent\chapnumfont\thechapter}
  \renewcommand{\afterchapternum}{\par\vskip\midchapskip}
  \renewcommand{\printchaptertitle}[1]{\chaptitlefont ##1\par
    \vskip 10pt {\color{accent}\rule{\textwidth}{2pt}}}
}
\makeatother
\chapterstyle{c64}

\setsecheadstyle{\sffamily\bfseries\fontsize{14.5}{18}\selectfont\color{c64ink}\raggedright}
\setsubsecheadstyle{\sffamily\bfseries\fontsize{11.8}{15}\selectfont\color{accent}\raggedright}
\setsubsubsecheadstyle{\sffamily\bfseries\fontsize{10.6}{14}\selectfont\color{c64muted}\raggedright}
\setbeforesecskip{-1.5\onelineskip}
\setaftersecskip{0.4\onelineskip}
\setsecnumdepth{subsection}
\maxsecnumdepth{subsection}
\renewcommand{\thesection}{\thechapter.\arabic{section}}

% -------------------------------------------------------- running head
\makepagestyle{c64page}
\makeevenhead{c64page}{\sffamily\small\color{c64muted}\thepage}{}{\sffamily\small\color{c64muted}\leftmark}
\makeoddhead{c64page}{\sffamily\small\color{c64muted}\rightmark}{}{\sffamily\small\color{c64muted}\thepage}
\makeheadrule{c64page}{\textwidth}{0.4pt}
\makeevenfoot{c64page}{}{}{}
\makeoddfoot{c64page}{}{}{}
\makepsmarks{c64page}{%
  \createmark{chapter}{both}{shownumber}{}{\quad}
  \createmark{section}{right}{shownumber}{}{\quad}}
\pagestyle{c64page}
\copypagestyle{chapter}{empty}

% --------------------------------------------------------------- lists
\setlist[itemize]{leftmargin=1.1em,itemsep=0.15\baselineskip,topsep=0.2\baselineskip,parsep=0pt}
\setlist[enumerate]{leftmargin=1.5em,itemsep=0.15\baselineskip,topsep=0.2\baselineskip,parsep=0pt}
\renewcommand{\labelitemi}{\textcolor{accent}{\rule[0.25ex]{0.42em}{0.42em}}}

% ------------------------------------------------------------- figures
% Screenshots are phone captures: tall, and unreadable if stretched to
% the measure. \screenshot picks a width from the shape of the file, and
% the frame keeps a light-mode capture from bleeding into the page.
\usepackage[font={sf,small,color=c64muted},labelfont={color=accent,bf},labelsep=period,skip=7pt]{caption}
\newcommand{\screenshot}[2]{%
  \begingroup
  \setlength{\fboxsep}{0pt}\setlength{\fboxrule}{0.6pt}%
  \textcolor{c64rule}{\fbox{\includegraphics[width=#1]{#2}}}%
  \endgroup}

% -------------------------------------------------------------- tables
% Sans, a touch smaller than the body, with a tinted header band. No
% vertical rules: they fence the reader in and buy nothing.
\newcommand{\tablehead}[1]{\sffamily\bfseries\small\color{c64ink}#1}
\renewcommand{\arraystretch}{1.28}
\setlength{\tabcolsep}{7pt}

% ------------------------------------------------------------ callouts
\newtcolorbox{manualnote}[1][]{%
  enhanced, breakable, sharp corners, boxrule=0pt,
  leftrule=2.4pt, colframe=accent, colback=accent!5!c64paper,
  left=10pt, right=10pt, top=8pt, bottom=8pt,
  fontupper=\small, #1}

% --------------------------------------------------------------- index
% A letter heading in the chapter accent, so the index reads as part of the
% same book rather than as an appendix bolted on.
\newcommand{\indexletter}[1]{%
  \par\vspace{11pt}%
  {\sffamily\bfseries\normalsize\color{accent}#1}%
  \par\vspace{2pt}{\color{c64rule}\rule{\linewidth}{0.4pt}}\par\vspace{4pt}}
\renewcommand{\preindexhook}{\small}

% ---------------------------------------------------------------- misc
\newcommand{\tightpar}{\vspace{-0.35\baselineskip}}
\hyphenpenalty=200
\tolerance=1600
\emergencystretch=1.2em
